% ============================================================
% 分片：s03a-model  对应原文 PDF 第 16--19 页（印刷页码 16--19）
% 负责章节：§3 THE C-Y/L-G CORRESPONDENCE（引言）+ 3.1 The Model
% 译者：t04-s3a   审校状态：初稿
% 审校：t02-s2a 已完成逐式核对（修正 5 处，另增 1 条存疑批注）
% ============================================================

% ==== 原文 PDF 第 16 页（印刷页 16）====
% [本页 §3 标题之前的两段属 s02b 分片；本片自 §3 标题起。]

\section{Calabi--Yau 与 Landau--Ginzburg 的对应}

\begin{keytakeaway}
\textbf{核心物理图像：}二维 $N=2$ 超对称的 $U(1)$ 规范线性 $\sigma$ 模型（GLSM）把
Calabi--Yau 非线性 $\sigma$ 模型与 Landau--Ginzburg 轨形统一成同一理论的两个相。
模型由 $n$ 个荷为 $1$ 的手征超场 $S_i$ 与一个荷为 $-n$ 的场 $P$ 构成，超势为
$W=P\,G(S_1,\dots,S_n)$，其中 $G$ 是 $n$ 次齐次横截多项式；FI 参数
（Fayet--Iliopoulos 参数）$r$ 的符号决定理论落在哪一个相。

\par\textbf{数学工具：}消去辅助场得到 $D$ 项与 $F$ 项势能；$D$ 平坦条件
$\sum_i\bar s_is_i=r$ 对 $U(1)$ 的辛商给出射影空间 $\mathbf{CP}^{n-1}$，
$F$ 平坦条件 $G=0$ 再把它截为 $n$ 次超曲面 $X$；$K_X=\mathcal{O}(k-n)|_X$
给出 Calabi--Yau 判据；$r<0$ 时的残余 $\mathbb{Z}_n$ 规范群即轨形群。

\par\textbf{主要结论：}$r\gg0$ 的低能理论是靶空间为 $X\subset\mathbf{CP}^{n-1}$、凯勒类正比于 $r$
的非线性 $\sigma$ 模型；$r\ll0$ 时得到在原点有 $n$ 阶退化临界点的 Landau--Ginzburg 轨形。
二者是同一个规范线性 $\sigma$ 模型的两个相。
\end{keytakeaway}

现在转向本文的问题：对 Calabi--Yau 模型与 Landau--Ginzburg 模型之间的对应关系作出一些阐述。
为此，我们只考虑最简单的情形，即规范群就是 $G=U(1)$，从而只有一个矢量超场 $V$。
各种推广留待 §4 讨论。

所幸，在 §2 的全部公式中，我们主要只需要少数几个比较简单的公式。场 $D$ 与 $F_i$
作为“辅助场”出现在拉格朗日量 (2.17) 中，不带动能项。可以利用它们的运动方程把它们解出，
得到
\begin{equation}
\tag{3.1}
\begin{aligned}
D &= -e^2\Bigl(\sum_i Q_i|\phi_i|^2-r\Bigr)\\
F_i &= \frac{\partial W}{\partial\phi_i}.
\end{aligned}
\end{equation}
% [Tutor Note: (3.1) 的两式就是作用量对辅助场的变分方程：因 (2.17) 中 $D$ 与 $F_i$ 只以代数方式出现（无导数项），这两个方程是代数方程，可直接解出并代回；把 (3.1) 代回 (3.2) 中的 $\frac{1}{2e^2}D^2$ 与 $\sum_i|F_i|^2$，即得通常的 $D$ 项势与 $F$ 项势。注意 $D$ 的定义式自带因子 $e^2$，代回后 $U$ 中的 $D$ 项为 $\frac{e^2}{2}\bigl(\sum_iQ_i|\phi_i|^2-r\bigr)^2$，故 $r$ 直接控制 $r>0$ 时 $s_i$ 的典型量级 $|s_i|^2\sim r/n$。]

于是，动力学标量场 $\phi_i,\sigma$ 的势能便可写成
\begin{equation}
U(\phi_i,\sigma)=\frac{1}{2e^2}D^2+\sum_i|F_i|^2+2\bar\sigma\sigma\sum_i Q_i^2|\phi_i|^2.
\tag{3.2}
\end{equation}

% ==== 原文 PDF 第 17 页（印刷页 17）====
\label{op:17}

\subsection{模型}

现在说明下面将要关注的情形。我们把手征超场 $\Phi_i$ 取为 $n$ 个荷为 $1$ 的场 $S_i$，
再加一个荷为 $-n$ 的场 $P$。这显然保证了反常抵消与 $R$ 不变性所需的条件 (2.34)。
小写字母 $s_i$ 与 $p$ 将表示超多重态 $S_i$ 与 $P$ 中的玻色场。
我们把超势取为规范不变的函数
\begin{equation}
W=P\cdot G(S_1,\dots,S_n),
\tag{3.3}
\end{equation}
其中 $G$ 是 $n$ 次齐次多项式。注意这个 $W$ 是拟齐次的；可以把 (2.35) 中的 $k_i$ 取为：$P$ 取 $-1$、$S_i$ 取 $0$；
或者（加上规范生成元的某个倍数之后）$P$ 取 $0$、$S_i$ 取 $-1/n$。
本文后面所考虑的超势总以类似显然的方式拟齐次；在后面的例子中不再详细说明这一点。
% [Tutor Note: 荷的取值 $\{Q_{S_i}=1,\ Q_P=-n\}$ 给出 $\sum_iQ_i=n-n=0$，这正是 (2.34) 的无反常条件；配合 $k_P=-1$、$k_{S_i}=0$ 的 $R$ 荷，$W=P\,G$ 满足 $W\to e^{-i\alpha}W$，故 $W$ 拟齐次。二者的共同来源是：$P$ 的荷（因而 $G$ 的次数）被无反常条件锁死为 $n$，这正是下一段 $k=n$ 的 Calabi--Yau 条件的伏笔。]

我们希望把 $G$ 选为“横截的”，意即方程
\begin{equation}
0=\frac{\partial G}{\partial S_1}=\dots=\frac{\partial G}{\partial S_n}
\tag{3.4}
\end{equation}
除 $S_i=0$ 之外没有公共根。这一性质保证：如果把 $S_i$ 看作齐次坐标，
则由方程 $G=0$ 在 $\mathbf{CP}^{n-1}$ 中定义的超曲面 $X$ 是光滑的。
一般的齐次多项式都具有这一性质。在我们后面的每一个模型中都会出现一个类似的横截性条件。
% [Tutor Note: 横截性 $\Rightarrow$ $X$ 光滑：由 Euler 关系 $\sum_iS_i\,\partial G/\partial S_i=nG$，在 $X$ 上 $G=0$；若 $X$ 上某点同时满足 $\partial G/\partial S_i=0$（$i=1,\dots,n$），该点即为 $X$ 的奇点。反之，梯度在 $X$ 上处处不为零，由隐函数定理知 $X$ 光滑。注意 (3.4) 用大写超场 $S_i$，(3.5) 起改用玻色分量 $s_i$，核对时勿混。]

对刚刚引入的模型，玻色势 (3.2) 为
\begin{equation}
\tag{3.5}
U=|G(s_i)|^2+|p|^2\sum_i\left|\frac{\partial G}{\partial s_i}\right|^2
+\frac{1}{2e^2}D^2+2|\sigma|^2\Bigl(\sum_i|s_i|^2+n^2|p|^2\Bigr)
\end{equation}
其中
\begin{equation}
D=-e^2\Bigl(\sum_i\bar s_is_i-n\bar pp-r\Bigr).
\tag{3.6}
\end{equation}
% [Tutor Note: 由 (3.2) 到 (3.5) 的核对链：$F_p=\partial W/\partial p=G$ 给出 $|G|^2$；$F_{s_i}=\partial W/\partial s_i=p\,\partial G/\partial s_i$ 给出 $|p|^2\sum_i|\partial G/\partial s_i|^2$；$D$ 项照抄；最后一项由 $\sum_iQ_i^2|\phi_i|^2=\sum_i|s_i|^2+n^2|p|^2$ 得到。$U=0$ 要求 $D=0$、$G=0$、$p\,\partial_iG=0$ 与 $\sigma=0$ 同时成立，这四组条件分别对应下文的真空方程。]

现在讨论低能物理随 $r$ 的不同取值有何变化。首先取 $r\gg0$。
此时要使 $D=0$，就要求 $s_i$ 不能全为零。

% ==== 原文 PDF 第 18 页（印刷页 18）====
\label{op:18}

既然如此，欲使势中 $|p|^2\sum_i|\partial_iG|^2$ 这一项消失，就必须（鉴于 $G$ 的横截性）有 $p=0$。
于是 $D$ 的消失恰好给出
\begin{equation}
\sum_i\bar s_is_i=r.
\tag{3.7}
\end{equation}
把 (3.7) 的解空间对规范群 $U(1)$ 取商，恰好得到复射影空间 $\mathbf{CP}^{n-1}$ 的一个副本，
其凯勒类正比于 $r$。最后还须令 $G=0$，以保证势的公式 (3.5) 中第一项消失；
并令 $\sigma=0$，以保证最后一项消失。
% [Tutor Note: 几何解释：方程 (3.7) 定义球面 $S^{2n-1}=\{\sum_i|s_i|^2=r\}$（$r>0$），$U(1)$ 在其上自由作用，辛商（symplectic quotient，矩映射水平为 $r$）恰为 $\mathbf{CP}^{n-1}$；$r$ 因而就是商空间上 Fubini--Study 凯勒形式的上同调类。注意 $r<0$ 时 (3.7) 无解，$D$ 平坦条件改由 $p\neq0$ 承担（见下文）。若计入 $\theta$ 角，$r$ 与 $\theta$ 合起来构成复化凯勒类。]

于是，经典真空的空间同构于由 $G=0$ 定义的超曲面 $X\subset\mathbf{CP}^{n-1}$。
众所周知，$\mathbf{CP}^{n-1}$ 中的 $k$ 次光滑超曲面是 Calabi--Yau 流形，当且仅当 $k=n$。
我们并非偶然撞上了 Calabi--Yau 条件。我们把 $P$ 取为荷 $-n$、从而把 $G$ 取为 $n$ 次，
正是为了保证无反常 $R$ 不变性所需的条件 (2.34)。底层模型的无反常 $R$ 不变性
保证了有效低能 $\sigma$ 模型的这种不变性；而对 $\sigma$ 模型而言，
无反常的 $R$ 不变性等价于 Calabi--Yau 条件。
% [Tutor Note: “众所周知”的判据可由伴随公式补足：$K_{\mathbf{CP}^{n-1}}=\mathcal{O}(-n)$，故 $k$ 次超曲面 $X$ 的典范丛为 $K_X=\mathcal{O}(k-n)|_X$，$K_X$ 平凡等价于 $c_1(X)=0$，即 $k=n$。这也说明为何无反常 $R$ 不变性等价于 Calabi--Yau 条件：二维 $\sigma$ 模型中 $U(1)_R$ 的 't Hooft 反常正比于 $c_1(X)$。]

除沿 $X$ 的切向振动之外，其余模式在树级都有质量。（规范场 $v$ 通过希格斯机制获得质量；
$p$、$\sigma$ 以及 $s_i$ 的非切向模式的质量可在 (3.4) 中看出。）
因此低能理论是以 $X$ 为靶空间、凯勒类正比于 $r$ 的 $\sigma$ 模型。
% [存疑: 原文此处称质量“可见于 (3.4)”，但 (3.4) 是横截性条件；$p$、$\sigma$ 与 $s_i$ 非切向模式的质量项实际出现在势 (3.5) 的 $|G|^2$ 与 $2|\sigma|^2(\sum_i|s_i|^2+n^2|p|^2)$ 等项中，疑为 (3.5) 之误。照抄原样，待核。]

现在考虑 $r\ll0$ 的情形。此时 $D$ 的消失要求 $p\neq0$。
于是 $|p|^2\sum_i|\partial_iG|^2$ 的消失（在 $G$ 横截的前提下）要求所有 $s_i=0$。
在这种情况下，$p$ 的模必须是 $|p|=\sqrt{-r/n}$。通过一个规范变换，可以把 $p$ 的辐角固定为零。
因此，该理论在规范变换意义下有唯一的经典真空。在这个真空附近展开时，
$s_i$ 是无质量的（对 $n\ge3$）。这些无质量场由某个有效超势支配，
该超势可以通过积掉有质量场 $p$ 来确定；在此情形下，积掉 $p$ 无非就是把 $p$ 取为它的期望值。
% [Tutor Note: $n\ge3$ 时 $s_i$ 无质量：$\partial G/\partial s_i$ 在原点的最低阶为 $n-1$ 次，故 $|p|^2|\partial_iG|^2$ 给出的质量项形如 $|s|^{2n-4}$，$n\ge3$ 时不存在二次项；$n=2$ 时该项恰给出 $s_i$ 的二次质量。这正是原文把 $n\ge3$ 单列出来的原因。]

% ==== 原文 PDF 第 19 页（印刷页 19）====
\label{op:19}

于是低能理论的有效超势为 $\widetilde{W}=\sqrt{-r}\cdot W(s_i)$。因子 $\sqrt{-r}$ 无关紧要，
因为它可以通过对 $s_i$ 重新标度而吸收掉。这个有效超势（对 $n\ge3$）在原点有退化临界点，
在原点处直到 $n$ 阶都为零。一个具有唯一经典真空态、且由带有退化临界点的超势支配的理论，
通常称为 Landau--Ginzburg 理论。
% [存疑: 第 18 页给出 $|p|=\sqrt{-r/n}$，而本页两处均写作 $\sqrt{-r}$（$p=\sqrt{-r}$ 与 $\widetilde{W}=\sqrt{-r}\cdot W(s_i)$）；两者相差常数因子 $1/\sqrt n$，$n\neq1$ 时不能同时成立，疑为原文代入 $p$ 的期望值时略去了 $1/\sqrt n$。该常数可被 $s_i$ 的重标度吸收，故不影响后文结论；照抄原样，待核。]
% [存疑: 按 (3.3)，$W=P\cdot G(S_1,\dots,S_n)$；积掉 $p$（即取 $p$ 为其期望值 $\sqrt{-r}$）后应得 $\widetilde{W}=\sqrt{-r}\,G(s_i)$。原文写作 $W(s_i)$，疑为 $G(s_i)$ 之误，也可能是“$p$ 取期望值后的 $W$”之简写。照抄原样，待核。]

在此情形下，我们实际上得到的是 Landau--Ginzburg 轨形，即 Landau--Ginzburg 理论的轨形，
理由如下。$p$ 的真空期望值并未完全破坏规范不变性；它把 $U(1)$ 破缺到一个
$\mathbb{Z}_n$ 子群 $E$，后者按 $s_i\to\zeta s_i$ 作用，其中 $\zeta$ 是 $1$ 的一个 $n$ 次根。
残余的规范不变性意味着我们实际得到的是 Landau--Ginzburg 理论的 $\mathbb{Z}_n$ 轨形。
这可由一个普遍而初等、尽管也许尚未被普遍承认的事实推出：轨形（按该词在量子场论中的通常用法）
等价于具有有限规范群的理论。
% [Tutor Note: 残余群确为 $\mathbb{Z}_n$：在规范变换 $s_i\to e^{i\lambda}s_i$、$p\to e^{-in\lambda}p$ 之下，保持 $\langle p\rangle\neq0$ 要求 $e^{-in\lambda}=1$，即 $e^{i\lambda}=\zeta$ 为 $n$ 次单位根，故 $E\cong\mathbb{Z}_n$。而“对单值性求和 = 对 $E$-twist 求和”正是“轨形 = 有限规范群理论”的实质：$\Sigma$ 上的 $E$ 单值性即各沟道中的 $E$-twist 扇区。]

事实上，在黎曼面 $\Sigma$ 上计算路径积分时，我们不必在作用量的绝对极小处
$p=\sqrt{-r}$、$s_i=v=0$ 及其全部规范变换附近展开，而可以在“只在局部才与
$p=\sqrt{-r}$、$s_i=v=0$ 规范等价”的位形附近展开。在这种位形中，规范场 $v$ 可以带有单值性。
由于单值性必须保持 $p$ 不变，它们取值于 $n$ 次单位根构成的群 $E$。路径积分的低能近似构造
包含对一切可能的 $E$ 值单值性求和。但这个对单值性的求和，恰恰就是构造 $E$ 轨形量子场论时
在所有沟道中对 $E$-twist 求和的通常操作。因此，一般而言，轨形可以看作规范群为有限群的
规范场论的一个特例。
% [Tutor Note: 原文以 $v$ 记 $U(1)$ 规范场（$v_\mu$ 为规范势），而非矢量多重态中的复标量 $\sigma$，故此处“$s_i=v=0$”应与 (3.5) 中要求的 $\sigma=0$ 分开理解。“$E$ 值单值性”指绕 $\Sigma$ 的非平凡闭路一周后，位形只回升到相差一个 $E$ 变换者；$\Sigma$ 的每个手柄可独立选择一次，对各扇区求和即得轨形的 twisted 构造。]

当然，无论 $r\gg0$ 还是 $r\ll0$，为得到无质量场的有效理论而对有质量场做积分，
都不能纯粹按经典方式进行；必须考虑对上面得到的 $\sigma$ 模型或 Landau--Ginzburg 轨形的
量子修正。我们将在 §3.2 中更仔细地考察这一点；这里只需指出% <-- 行末 % 防止跨文件拼接产生空格（s03b 紧接续译）
% [跨页接续: 原文第 19 页/第 20 页；本句在下一页续完，其余部分归 s03b 分片。]
